
\documentclass[12pt]{article}
\usepackage{amssymb}
\usepackage{geometry}
\usepackage{graphicx}
\usepackage[colorlinks=true, allcolors=blue]{hyperref}
\usepackage{cite}
\usepackage{wrapfig}
\usepackage{booktabs}
\usepackage{soul}
\usepackage[most]{tcolorbox}
\usepackage{amsmath}
\usepackage{lipsum}
\usepackage{caption}
\usepackage{ulem}
\usepackage{xcolor}
\usepackage{tikz}
\usetikzlibrary{arrows.meta, positioning, shapes}
\usepackage{sectsty}

%--------- Define Colors ---------
\definecolor{titleColor}{HTML}{800000}        % Maroon
\definecolor{sectionColor}{HTML}{4682B4}      % SteelBlue
\definecolor{subsectionColor}{HTML}{6A5ACD}   % SlateBlue
\definecolor{textHighlight}{HTML}{DAA520}     % Goldenrod
\definecolor{backgroundColor}{HTML}{F0F8FF}   % AliceBlue
\definecolor{emphasisColor}{RGB}{200, 50, 50} % ForestGreen

\sectionfont{\color{sectionColor}}       % Color for \section
\subsectionfont{\color{subsectionColor}} % Color for \subsection

%--------- Custom Highlight Commands ---------
\definecolor{lightblue}{RGB}{230, 240, 255}
\definecolor{lightgreen}{RGB}{220, 255, 220}
\definecolor{lightred}{RGB}{255, 230, 230}
\definecolor{lightyellow}{RGB}{255, 255, 200}
\definecolor{darkblue}{RGB}{0, 51, 102}

\newcommand{\bluehighlight}[1]{\colorbox{lightblue}{#1}}
\newcommand{\greenhighlight}[1]{\colorbox{lightgreen}{#1}}
\newcommand{\redhighlight}[1]{\colorbox{lightred}{#1}}
\newcommand{\yellowhighlight}[1]{\colorbox{lightyellow}{#1}}
\newcommand{\blue}[1]{\textcolor{darkblue}{#1}}
\newcommand{\myemphasis}[1]{\textcolor{emphasisColor}{\textbf{#1}}}

%--------- Color Blocks ---------
\newtcolorbox{blockA}{
  breakable,
  colback=blue!5,
  colframe=blue!40!black,
  enhanced,
  boxrule=0.7pt,
  arc=2mm,
  left=4mm,right=4mm,top=2mm,bottom=2mm
}

\newtcolorbox{blockB}{
  breakable,
  colback=red!5,
  colframe=red!40!black,
  enhanced,
  boxrule=0.7pt,
  arc=2mm,
  left=4mm,right=4mm,top=2mm,bottom=2mm
}

\newtcolorbox{blockC}{
  breakable,
  colback=green!5,
  colframe=green!40!black,
  enhanced,
  boxrule=0.7pt,
  arc=2mm,
  left=4mm,right=4mm,top=2mm,bottom=2mm
}

\newtcolorbox{blockD}{
  breakable,
  colback=purple!5,
  colframe=purple!40!black,
  enhanced,
  boxrule=0.7pt,
  arc=2mm,
  left=4mm,right=4mm,top=2mm,bottom=2mm
}

%--------- Title ---------
\title{\textcolor{titleColor}{\textbf{\Huge Pharmacy V13 integrated  management-intelligence platform }}}
\author{B. M. Osman\thanks{Al Nahada University, Cairo \\ \texttt{babikerosman@yahoo.com}}}
\date{\today}

\begin{document}
\maketitle


\begin{abstract}

\textbf{Study context and scope.}
Pharmacy V13 is developed in the context of pharmacy operations where
purchasing, sales, inventory, suppliers, expiry, and financial transactions
generate substantial operational information but do not, by themselves,
provide integrated management intelligence. The study addresses the need to
transform these routine operational records into structured, interpretable,
and evidence-traceable decision support. The scope encompasses the design and
integration of an operational pharmacy information system with analytical and
management-intelligence capabilities, with particular emphasis on supply
chain analysis, inventory management, demand understanding, risk assessment,
and model evaluation. The system is intended to support management rather than
to replace professional judgement or authorise automated decisions.

Pharmacy V13 is therefore conceived as an integrated management intelligence
system that transforms routine pharmacy operations into structured,
evidence-based decision support. Rather than functioning solely as a system
for recording medicines, purchases, sales, inventory, suppliers, and financial
transactions, V13 establishes a progressive analytical architecture linking
operational evidence to information, analytical understanding, management
intelligence, and accountable decision-making.

The system incorporates a dedicated Supply Chain Analysis and Intelligence
Layer, supported by the REEM intelligence engine, to integrate observed
purchasing, sales, inventory, supplier, expiry, and demand evidence. Detailed
computational outputs from analytical models are transformed into management
summaries, visual analytics, risk and exception indicators, model-readiness
assessments, and evidence-traceable management signals. The analytical
portfolio includes demand forecasting, stock coverage, ABC analysis, expiry
risk, reorder point, safety stock, economic order quantity (EOQ), and Monte
Carlo uncertainty analysis.

A central principle of the architecture is the distinction between observed
evidence, derived indicators, model outputs, and decision intelligence. Model
execution alone is therefore not considered sufficient for operational use.
Models are assessed according to their analytical purpose and evidence
requirements. Continuous forecasting models are evaluated through forecast
error and residual analysis, optimisation models through parameter and
sensitivity analysis, while appropriate binary classification models may be
evaluated using confusion matrices, precision, recall, sensitivity,
specificity, F1 score, calibration, and receiver operating characteristic
(ROC) analysis with area under the curve (AUROC), where adequate labelled
historical outcomes exist.

The architecture adopts a progressive model-readiness framework ranging from
\textit{Available} and \textit{Partial} to \textit{Ready for Extension} and
\textit{Operationally Validated}. This prevents mathematically completed
calculations from being interpreted automatically as validated procurement
recommendations. Instead, Pharmacy V13 positions analytical models as
transparent instruments for management understanding, risk assessment, and
decision support.

The study does not claim that every analytical model is operationally
validated at the point of implementation. The maturity of each model depends
on the availability, quality, duration, and completeness of historical
evidence, together with the availability of appropriate model parameters and
labelled outcomes where required. Consequently, exploratory and descriptive
analyses may be used where evidence is currently limited, while predictive,
classification, simulation, and optimisation capabilities are progressively
validated as the evidence base matures.

The overall objective is to establish a scientifically defensible and
evidence-traceable pathway from observed pharmacy transactions to reliable
historical knowledge, validated analytical models, risk-aware prediction,
responsible optimisation, and accountable management action. Within this
scope, artificial intelligence, statistical modelling, simulation, and
decision-support methods are intended to inform management; they do not
replace human judgement, professional responsibility, or organisational
governance.

\end{abstract}
\end{document}
